\begin{ex}
Số tự nhiên bé nhất là số nào?
\shortans{0}
\loigiai{
Số tự nhiên bé nhất là 0.
}
\end{ex}

\begin{ex}
Có số tự nhiên lớn nhất không? (Có ghi 1, Không ghi 0).
\shortans{0}
\loigiai{
Không có số tự nhiên lớn nhất vì dãy số tự nhiên kéo dài vô tận.
}
\end{ex}

\begin{ex}
Hai số tự nhiên liên tiếp hơn kém nhau bao nhiêu đơn vị?
\shortans{1}
\loigiai{
Hơn kém nhau 1 đơn vị.
}
\end{ex}

\begin{ex}
Hai số chẵn liên tiếp hơn kém nhau bao nhiêu đơn vị?
\shortans{2}
\loigiai{
Hơn kém nhau 2 đơn vị.
}
\end{ex}

\begin{ex}
Hai số lẻ liên tiếp hơn kém nhau bao nhiêu đơn vị?
\shortans{2}
\loigiai{
Hơn kém nhau 2 đơn vị.
}
\end{ex}

\begin{ex}
Số liền sau của số tự nhiên $n$ là biểu thức nào? (Nếu $n=10$ thì số liền sau là bao nhiêu?).
\shortans{11}
\loigiai{
$10 + 1 = 11$.
}
\end{ex}

\begin{ex}
Số liền trước của số tự nhiên $n$ ($n > 0$) với $n=20$ là bao nhiêu?
\shortans{19}
\loigiai{
$20 - 1 = 19$.
}
\end{ex}

\begin{ex}
Dãy số tự nhiên bắt đầu từ số nào?
\shortans{0}
\loigiai{
Bắt đầu từ số 0.
}
\end{ex}

\begin{ex}
Trong dãy số $0, 1, 2, 3, 4, \dots$, số thứ 10 là số nào? (Số thứ nhất là 0).
\shortans{9}
\loigiai{
Số thứ nhất là 0, thứ 10 là 9.
}
\end{ex}

\begin{ex}
Trong dãy số $1, 2, 3, 4, \dots$, số thứ 50 là số nào?
\shortans{50}
\loigiai{
Số thứ 50 là 50.
}
\end{ex}

\begin{ex}
Dãy số $2, 4, 6, 8, \dots$ là dãy số gì? (Dãy số chẵn ghi 0, dãy số lẻ ghi 1).
\shortans{0}
\loigiai{
Dãy số chẵn (0).
}
\end{ex}

\begin{ex}
Dãy số $1, 3, 5, 7, \dots$ là dãy số gì? (Dãy số chẵn ghi 0, dãy số lẻ ghi 1).
\shortans{1}
\loigiai{
Dãy số lẻ (1).
}
\end{ex}

\begin{ex}
Số hạng tiếp theo của dãy số $2, 4, 6, 8, \dots$ là số nào?
\shortans{10}
\loigiai{
Số tiếp theo là 10.
}
\end{ex}

\begin{ex}
Số hạng tiếp theo của dãy số $1, 3, 5, 7, \dots$ là số nào?
\shortans{9}
\loigiai{
Số tiếp theo là 9.
}
\end{ex}

\begin{ex}
Số hạng tiếp theo của dãy số $5, 10, 15, 20, \dots$ là số nào?
\shortans{25}
\loigiai{
Số tiếp theo là 25.
}
\end{ex}

\begin{ex}
Số hạng tiếp theo của dãy số $1, 4, 7, 10, \dots$ là số nào?
\shortans{13}
\loigiai{
Mỗi số hơn kém nhau 3 đơn vị: $10 + 3 = 13$.
}
\end{ex}

\begin{ex}
Số hạng tiếp theo của dãy số $1, 2, 4, 8, 16, \dots$ là số nào?
\shortans{32}
\loigiai{
Gấp đôi số trước: $16 \times 2 = 32$.
}
\end{ex}

\begin{ex}
Có bao nhiêu số tự nhiên có một chữ số?
\shortans{10}
\loigiai{
Từ 0 đến 9 có 10 số.
}
\end{ex}

\begin{ex}
Có bao nhiêu số tự nhiên có hai chữ số?
\shortans{90}
\loigiai{
Từ 10 đến 99 có $(99 - 10) + 1 = 90$ số.
}
\end{ex}

\begin{ex}
Có bao nhiêu số tự nhiên có ba chữ số?
\shortans{900}
\loigiai{
Từ 100 đến 999 có 900 số.
}
\end{ex}

\begin{ex}
Có bao nhiêu số tự nhiên từ 1 đến 100?
\shortans{100}
\loigiai{
$100 - 1 + 1 = 100$ số.
}
\end{ex}

\begin{ex}
Có bao nhiêu số chẵn từ 1 đến 100?
\shortans{50}
\loigiai{
$100 : 2 = 50$ số chẵn.
}
\end{ex}

\begin{ex}
Có bao nhiêu số lẻ từ 1 đến 100?
\shortans{50}
\loigiai{
$100 : 2 = 50$ số lẻ.
}
\end{ex}

\begin{ex}
Số 0 có phải là số tự nhiên không? (Phải ghi 1, Không ghi 0).
\shortans{1}
\loigiai{
0 là số tự nhiên.
}
\end{ex}

\begin{ex}
Trong các số tự nhiên từ 10 đến 20 có bao nhiêu số?
\shortans{11}
\loigiai{
$20 - 10 + 1 = 11$ số.
}
\end{ex}
